\chapter{Metode}

\section{Forskningsdesign}

Studien gjennomføres som en evaluativ brukbarhetsstudie av en mobilapplikasjon for datainnsamling. Målet er å identifisere friksjon og feil knyttet til systemtilstand og vurdere konsekvenser for datakonsistens.

Designet kombinerer:

\begin{itemize}
\item Scenario-basert brukertesting
\item Think-aloud-protokoll
\item Semistrukturert intervju
\end{itemize}

Tilnærmingen gjør det mulig å dokumentere både observerbar atferd og brukerens forståelse og forventninger.

\section{Deltakere og rekruttering}

Utvalget er strategisk for å belyse forskjeller mellom tekniske og ikke-tekniske brukere. Planlagt utvalg inkluderer:

\begin{itemize}
\item Ikke-tekniske brukere fra avfalls-/transportkontekst
\item Tekniske brukere med systemforståelse
\end{itemize}

Rekruttering skjer etter godkjenning fra Sikt.

\section{Testoppsett og datainnsamling}

Scenario-baserte oppgaver er konstruert for å fremprovosere kritiske tilstandsoverganger i applikasjonen (persistens, navigasjon, innsending og reset).

\subsection{Think-aloud}

Deltakere verbaliserer tanker under oppgaveløsning, noe som gir innsikt i mental modell og forventninger. Metoden er forankret i protokollanalyse \cite{ericsson1993}.

\subsection{Semistrukturert intervju}

Etter test gjennomføres intervju om forståelse av lagring, sending og reset. Dette utdyper observasjoner.

\section{Kvantitativ evaluering: to-trinns spørreundersøkelser}

\subsection{Metodisk tilnærming og datainnsamling}
Som en del av den empiriske evalueringen av dette prosjektet ble det gjennomført to anonyme spørreundersøkelser med identiske spørsmål basert på en Likert-skala fra 1 til 7. Disse spørsmålene ble utformet for å måle den faktiske brukeropplevelsen av StellAI-applikasjonen. Denne spørreundersøkelsen bygger på brukertesting gjennom scenarier og tenke-høyt-protokollen som ble beskrevet i kapittel fem, og danner et kvantitativt grunnlag for å fastslå om forbedringene som ble implementert, har hatt en merkbar og målbar effekt på brukervennligheten.

Den første runden av spørreundersøkelsen ble gjennomført i perioden fra 25. til 30. februar 2026, med deltakelse fra 31 personer. Deretter ble tilbakemeldingene som testpersonene ga i denne runden analysert og brukt som grunnlag for å gjennomføre målrettede endringer i applikasjonen. En tilsvarende spørreundersøkelse ble gjennomført i perioden fra 10. mars til 14. mars 2026, med deltakelse fra 21 personer. Fullføringsgraden for spørreundersøkelsen var 100\% i begge rundene. Det første spørsmålet (Q1) omhandlet deltakernes tekniske bakgrunn og ble brukt som en kontrollvariabel, mens spørsmål 2 til 10 (Q2–Q10) utgjorde evalueringssettet for applikasjonen og ga dermed grunnlag for sammenligning.

\subsection{Overordnet utvikling – fra 77,7\% til 86,2\%}
Gjennomsnittsvurderingen av applikasjonen for spørsmålene Q2 til Q10 økte fra 77,7\% til 86,2\%, noe som tilsvarer en forbedring på 8,5 prosentpoeng. Dette regnes som en positiv indikator på at utviklingen som ble gjort i applikasjonen – der brukernes tilbakemeldinger ble omgjort til konkrete designforbedringer og deretter evaluert på nytt – hadde en tydelig effekt. Samtidig ble gjennomsnittstiden som var nødvendig for å svare, redusert fra 4,0 til 3,2 minutter, noe som tyder på at applikasjonens brukervennlighet ble forbedret. Antallet fritekstkommentarer sank fra 26 til 17, og denne nedgangen gjenspeiler tydelig at sentrale svakheter i brukergrensesnittet ble forbedret, og ikke at brukernes engasjement ble redusert. Kontrollvariabelen Q1 (teknisk erfaring) forble stabil mellom de to rundene (61,3\% mot 59,2\%), noe som reduserer sannsynligheten for at forbedringene skyldes forskjeller i erfaringsnivået i utvalget etter testen. Dette er et viktig punkt for studiens validitet, ettersom et av hovedmålene med utformingen av StellAI er å støtte brukere med ulik teknisk bakgrunn – slik det er beskrevet i forskningsspørsmålet og i avsnitt 1.5 om brukervennlighet i oppgaveorienterte applikasjoner.

\subsection{Spørsmål for spørsmål: størst forbedringer}
\begin{itemize}
\item Q8 – Tydelige knapper, ikoner og tekster representerer den mest dramatiske enkeltendringen i hele datasette: Gjennomsnittsskåren økte fra 2,74 (39,2\%) til 6,38 (91,2\%), med en økning på 52,0 prosentpoeng. Dette er et direkte resultat av de hyppigst gjentatte kommentarene fra første runde, der 21 av 26 kommentarer på en eller annen måte pekte på fraværet av tekst under ikonene. Kommentarer som «ikonene mangler tekst, så det bør være tekst under dem», «sendeknappen og de andre ikonene (for å ta bilde eller laste opp) har ingen tekstlige forklaringer», og «det var ikoner som var irriterende fordi det ikke fantes tekst under dem», bekrefter at dette var det viktigste problemet knyttet til applikasjonens brukervennlighet. Her kan vi knytte dette resultatet direkte til det andre forskningsspørsmålet (RQ2) – hvordan utformingen av brukergrensesnittet påvirker brukerens forståelse av systemets tilstand – og til den teoretiske drøftingen i avsnitt 4.4 om Nielsens heuristiske prinsipper, der både «samsvar mellom system og virkelighet» og «gjenkjenning fremfor gjenkalling» anses som særlig viktige: Symbolene fremstod som abstrakte tegn uten kontekst som brukerne ikke klarte å gjenkjenne, noe som skapte en kognitiv barriere for effektiv oppgaveløsning.
\item Q7 – Tydelig fullført etter innsending steg fra 4,74 (67,7\%) til 6,33 (90,5\%), en forbedring på 22,7 prosentpoeng. Dette spørsmålet berører kjernen i det som i avsnitt 1.6 omtales som tydelige overganger mellom tilstander: Brukeren må forstå at registreringsprosessen er fullført, og at systemet er klart for en ny registrering. Andelen på 67,7\% i første runde bekrefter at dette var et reelt problem brukerne møtte, og usikkerhet rundt om innsendingen faktisk var fullført, bidro til gjentatte forsøk på innsending, som er en av mekanismene bak duplisering av data. Den tydelige forbedringen på 22,7 prosentpoeng etter endringene som ble gjort i applikasjonen, viser at bekreftelsesmeldingen og det visuelle elementet som signaliserer vellykket innsending, ble tydeligere, og dette bidro direkte til å forbedre kvaliteten på de innsendte dataene ved å redusere sannsynligheten for dupliserte registreringer.
\item Q3 – Bildeopplasting uten tap økte fra 5,03 (71,9\%) til 6,24 (89,1\%), en forbedring på 17,2 prosentpoeng. Gjentatte kommentarer i den første undersøkelsen, som «bildene slettes ikke etter innsending og vises igjen», «de forrige bildene vises etter innsending», og «bildene dupliseres når de sendes», forekom i sju av 26 kommentarer. Dette var en indikasjon på tap av datakvalitet og manglende konsistens, slik det ble drøftet i avsnitt 1.4 og 1.6: Når applikasjonen ikke tilbakestilles korrekt etter innsending, og tidligere bilder fortsatt vises i brukergrensesnittet, svekker dette brukerens forståelse av hva som faktisk er sendt, og hva som utgjør en ny registrering. Forbedringen i det tredje kvartilet etter redesignen peker på at tilbakestillingsmekanismen ble korrigert, og uttrykker dermed en direkte forbedring i datakvaliteten, som er applikasjonens grunnleggende mål.
\item Q9 – Trygghet via meldinger viste en moderat økning fra 5,74 (82,0\%) til 6,24 (89,1\%), en forbedring på 7,1 prosentpoeng. Dette spørsmålet omhandler også det teoretiske begrepet opplevd trygghet i databehandling, når brukeren sender bilder og metadata til skylagring, med en forventning om at disse dataene håndteres korrekt. Denne økningen tyder på at endringene som ble gjort i systemets kommunikasjon med brukeren – enten gjennom tydeligere bekreftelsesmeldinger eller mer informative statusmeldinger – bidro til å styrke tilliten.
\end{itemize}

\subsection{Dimensjoner med tilbakegang – kritisk analyse}
En objektiv analyse av dataene krever oppmerksomhet mot de dimensjonene som viste nedgang etter endringene.
\begin{itemize}
\item Q4 – Riktig kategori falt fra 6,52 (93,1\%) til 5,76 (82,3\%), en nedgang på 10,8 prosentpoeng. Denne nedgangen kan forklares direkte gjennom kommentarene som ble samlet inn i andre runde: «Valget for å endre kategori på bildet var skjult bak tastaturet, så jeg kunne ikke lese det jeg skrev. Når riktig kategori for bildet velges, endres den etter at man går tilbake til hovedsiden, og dette gjør det vanskelig å kontrollere om brukerens valg er korrekt.» Dette gjenspeiler en ny form for forvirring som oppsto som følge av forbedringene som ble gjort i funksjonen for kategorivalg. Denne endringen var ment å forbedre flyten i brukervennligheten, men førte til et nytt problem i utformingen av brukergrensesnittet ved at inntastingsfeltet ble skjult bak tastaturet, og bevegelsen i kategorien signaliserte usikkerhet snarere enn bekreftelse.
\item Q5 – Metadata gikk fra 6,87 (98,2\%) til 6,43 (91,8\%), en nedgang på 6,3 prosentpoeng. Det femte spørsmålet (Q5) var det høyest rangerte spørsmålet i hele første runde, og selv om vurderingen fortsatt er høy, kan vi si at den sank merkbart. Dette kan forklares med at endringene som ble gjort i skjermoppsettet rundt metadatafeltet – i tillegg til det nye problemet som oppsto i kategorivelgeren (jf.Q4) – førte til en viss grad av forvirring under registreringsprosessen. Det finnes også en annen mulig forklaring, nemlig at det femte spørsmålet (Q5) kan ha blitt påvirket av endringene i oppsettet og utformingen av brukergrensesnittet på siden for metadata. Selv om justeringene kan ha vært små, kan de ha ført til en reduksjon i den samlede skåren.
\item Q2 – forståelsen av rekkefølgen på trinnene falt svakt fra 6,06 (86,6\%) til 5,81 (83,0\%), altså med en nedgang på 3,6 prosentpoeng. Dette spørsmålet handler om arbeidsflyten i brukerens tenkning: bilde → kategori → metadata → innsending. Denne svake nedgangen i skår gjenspeiler at endringer i applikasjonens navigasjonsstruktur eller rekkefølgen på trinnene noen ganger kan skape en viss forvirring hos brukere som var vant til den tidligere arbeidsflyten.
\item Q10 – Opplevd sikkerhet falt fra 5,52 (78,8\%) til 5,29 (75,5\%), −3,3 prosentpoeng. Selv om dette er en liten endring, indikerer den en nedgang i en sammenheng der datainnsamling finner sted. Her bør brukeren føle høy grad av trygghet ved bruk av applikasjonen, for å sikre bruken av skylagring og trening av kunstig intelligens. Det er mulig at forbedringene som ble gjort i applikasjonens utseende, skapte en midlertidig følelse av fremmedgjøring som påvirket dette forholdet.
\end{itemize}

\subsection{Kvalitativ analyse av fritekstkommentarer}
I første runde hadde vi 26 kommentarer. Det ble lagt vekt på to hovedtemaer: uklarhet knyttet til knapper og ikoner (21 ganger), og gjentakelse av bilder etter innsending (7 ganger). Når 21 av 26 respondenter nevnte spørsmålet om ikoners tydelighet, gjenspeiler dette i seg selv en klar indikator, og dette begrunner hvordan dette temaet ble prioritert da det ble valgt som et sentralt endringspunkt. Når det gjelder de øvrige kommentarene som ble nevnt, for eksempel ønsket om innlogging via Google (4), funksjon for visningshistorikk og sletting (3), og tap av data (2), peker disse alle på mulige fremtidige forbedringer som det kan arbeides videre med, men de ble ikke behandlet direkte i iterasjonssyklusen som denne studien dekket.

I den andre fasen hadde vi 17 kommentarer. Endringene og forbedringene førte til at det tematiske bildet endret seg radikalt. Plasseringen av knappen for hovedsiden ble det hyppigst gjentatte temaet (6 kommentarer), og de fleste kommentarene var positive: «En flott app, enkel og elegant», «En utmerket og brukervennlig app», «En perfekt app», og «Utmerket når det gjelder tilbakemelding og et sterkt passord på 8 tegn. Flott app». Dette viser graden av tilfredshet blant en bred gruppe brukere. De kritiske kommentarene som ble mottatt, var derimot konsentrert rundt arbeidsmiljøet og mer presise forbedringer. Plasseringen av knappen for hovedsiden på høyre side utgjorde et hinder for høyrehendte brukere, noe som er en utfordring innen universell utforming, og klassifiseringsfeltet ble skjult bak tastaturet. Alle disse kommentarene er en indikator på at den kontinuerlige forbedringen har vært vellykket: Tilbakemeldingene ble mer detaljerte og spesifikke, noe som tyder på at brukerne ikke lenger var opptatt av mangler ved de grunnleggende funksjonene, men av å forbedre brukeropplevelsen.

\subsection{Kobling til forskningsspørsmål og teoretisk rammeverk}
Resultatene fra den komparative studien kan nå knyttes direkte til de fire forskningsspørsmålene som utgjør strukturen i denne avhandlingen.
\begin{itemize}
\item RQ1 - Her fremgår det hvordan ansatte med ulike nivåer av teknisk erfaring håndterer brukervennligheten i registreringsprosessen, gjennom forutsigbarhet, forventning om feil og fullføring av prosessen. Kontrollindikatoren Q1 (teknisk erfaring) forble stabil gjennom alle rundene, noe som gjorde det mulig å gjennomføre en sammenligning som ikke var påvirket av variasjoner i erfaring. Forbedringen i Q7 på +22,7 prosentpoeng og i Q8 på +52,0 prosentpoeng peker på en tydelig økning i prosessens forutsigbarhet og en fjerning av feilkilder knyttet til uklare symboler. Samtidig viser nedgangen i Q2 på -3,6 prosentpoeng og i Q4 på -10,8 prosentpoeng at den nye navigasjonsoppdateringen og kategorivelgeren førte til brukerforvirring som ikke tidligere var til stede, noe som gjør ytterligere forbedringer nødvendige.
\item RQ2 - I dette spørsmålet gjelder hvordan utformingen av brukergrensesnittet påvirker brukerens oppfatning av systemets tilstand. Endringen i Q8 representerer et direkte bevis: Fraværet av tekstlige etiketter førte til at brukerne ikke forsto knappenes funksjoner, og dermed til feiltolkning av systemets tilstand. På samme måte fremhever Q7 betydningen av en tydelig bekreftelse etter innsending, slik at brukeren kan forsikre seg om at systemet har gått fra tilstanden «ikke sendt» til tilstanden «sendt» – en overgang som i avsnitt 4.2 i det teoretiske rammeverket ble beskrevet som interaktiv innsending. Disse resultatene gjenspeiler samsvar med Nielsens og Hoas prinsipp om tydelig systemstatus og bekrefter de teoretiske forventningene som er presentert i teorikapitlet.
\item RQ3 - søker å dokumentere brukerfeil og fremvoksende datainkonsistenser, og i hvilken grad disse er knyttet til håndteringen av systemets tilstand som helhet og til brukerens mentale modell. Resultatene registrerte to typer feil som påvirket datakvaliteten: For det første førte knapper som var uklare for brukeren til feil handlinger (feil trykk, tvil om hvorvidt innsendingen var gjennomført, eller gjentatt innsending flere ganger), og for det andre førte manglende tilbakestilling i applikasjonen til at dupliserte bilder ble værende etter innsending. Disse to tilfellene viser hvordan unøyaktig formidling av systemets tilstand kan sette dataintegriteten i fare, og dette er nettopp spørsmålet som gjelder forholdet mellom brukervennlighet og datakvalitet.
\item RQ4 - Dette omhandler konkrete design- og implementeringsforbedringer som reduserer feil og styrker datakvaliteten. Responsen fra spørreundersøkelsen var tydelig: Tilføyelsen av tekstlige etiketter under symbolene (+52,0 prosentpoeng sammenlignet med Q8), tydeliggjøringen av innsendingbekreftelsen (+22,7 prosentpoeng sammenlignet med Q7), og korrigeringen av logikken for tilbakestilling av bilder (+17,2 prosentpoeng sammenlignet med Q3) utgjør de tre mest virkningsfulle dokumenterte tiltakene. Disse resultatene danner et empirisk grunnlag for designanbefalingene som vil bli presentert i diskusjonskapitlet (kapittel seks), særlig i avsnitt 6.4 som omhandler overgangen fra forskningsspørsmål til problemformulering.
\end{itemize}




\section{Analysemetode}

Data analyseres ved tematisk analyse \cite{braun2006}. Observasjoner, think-aloud og intervjudata kodes og grupperes i tema relatert til:

\begin{itemize}
\item Tilstandsuklarhet
\item Persistensproblemer
\item Resetproblemer
\item Mental model mismatch
\item Datainkonsistens
\end{itemize}

\section{Kvalitet i studien}

Troverdighet styrkes gjennom:

\begin{itemize}
\item Datatriangulering
\item Standardisert oppgavesett
\item Tydelig teori–analyse-kobling
\end{itemize}

\section{Forskerrolle og bias}

I denne studien har jeg en dobbeltrolle som utvikler av mobilapplikasjonssystemet og som gjennomfører brukertestene. Applikasjonen ble utviklet gjennom et samarbeid mellom Stellai AS og Universitetet i Agder (UiA), og løsningen ble designet og implementert i løpet av perioden for masteroppgaven. Denne rollen ga en dyp forståelse av hvordan systemets funksjoner fungerer og hvilken kontekst de inngår i, men innebærer også en risiko for skjevhet i utformingen av studien, gjennomføringen og tolkningen av resultatene.

\textbf{Design-bias.}
Siden forskeren selv har utviklet applikasjonen, kan forskeren være tilbøyelig til å vurdere designvalg mer positivt enn nødvendig, eller ubevisst favorisere den eksisterende applikasjonen i testoppsettet. For å forhindre denne muligheten ble oppgavene i brukertestene utformet for å fremkalle avgjørende tilstandsendringer i applikasjonen (persistens, navigasjon, innsending og tilbakestilling), uavhengig av forventet korrekt bruk. Resultatene ble analysert ved hjelp av forhåndsdefinerte indikatorer (tap av metadata, fiktive data, tilbakestillingsfeil og avbrudd i prosessen) for å sikre konsistent vurdering mellom alle deltakerne.

\textbf{Observatørbias.}
Når utvikleren selv fungerer som testleder, kan deltakerne endre atferden sin eller svarene sine, for eksempel ved å unngå kritikk eller forsøke å «gjøre det riktig». For å redusere denne effekten ble det presisert at det er systemet som testes, ikke deltakeren. I tillegg ble tenke-høyt-metoden brukt for å registrere spontane reaksjoner under oppgaveløsningen, og testlederen unngikk å forklare funksjoner eller gi veiledning.

\textbf{Reduksjon av tolkingsbias.}
Reduksjon av tolkningsskjevhet: Studien benytter datatriangulering ved å kombinere observasjon av oppgaveløsning, tenke-høyt-data og intervju­respons. Analysen ble gjennomført objektivt basert på konsepter fra teorikapittelet, slik at tolkningen bygger på systematiske mønstre i datamaterialet fremfor individuelle inntrykk.

